\scanpage{23}
\deleted{\unsure{Soit}} $K$ corps local d’inégales caractéristiques, corps résiduel $k$ de car. $p>0$ ($k$ parfait!), anneau d’entiers $V$, sous-$p$-anneau $W$ de corps résiduel $k$, corps des fractions $L$. On choisit une clôture algébrique $\bar K\deleted{\Omega}$ de $K$, son complété est $\Omega$.

\diagram{\begin{tikzcd}[ampersand replacement=\&,column sep=2pt,row sep=1pt]
L \& \subset \& K \& \subset \& \bar K \& \subseteq \& \Omega\\
\cup \& \& \cup \& \& \cup \& \& \cup\\
W \& \unsure{\in|\subset} \& V \& \subset \& \bar V \& \subset \& \widehat{\bar V}
\end{tikzcd}}

On pose
\[
\pi=\operatorname{Gal}(\bar K/K),\ \pi'=\operatorname{Gal}(\bar K/L)\,\ldots
\]
$\mathcal M_K=$ catégorie des \inserted{$\mathbb Q_p$-}faisceaux du type Hodge-Tate sur $K$ \inserted{($\otimes$-catégorie sur $\mathbb Q_p$)}

$\mathcal M_K^+=\underset{\longrightarrow}{\lim}\mathcal M^{(n)}$, ss-catégorie des objets ``effectifs'' i.e. avec \deleted{\unsure{la graduation|graduation}} $\operatorname{gr}(M_\Omega)$ à degrés positifs

$\mathcal M_K^{(\unsure{n|1})}=\bigl\{$ \deleted{\unsure{ss-}}catégorie pleine de $\mathcal M$ formée des $M$ avec $\operatorname{gr}M_\Omega$ à degrés $0$ et $1$

$\widetilde{\mathrm{BT}}_{\deleted{\unsure{K}}V}$ catégorie des BT à isog. près \uline{sur $V$}

Foncteur pl. fid (par Tate)
\[
\widetilde{\mathrm{BT}}_V\lhook\joinrel\longrightarrow\mathcal M_K^{(\unsure{1|n})}
\]
\margin{on identifie $\widetilde{\mathrm{BT}}_V$ à ss-catégorie pleine de $\mathcal M_K^{(1)}$}

$\mathcal M_{\deleted{\unsure{K+|K}}}^{\mathrm{BT}}\subset\mathcal M_K$ $\otimes$-\deleted{\unsure{ss-}}catégorie Tannakienne engendrée par $\mathcal M_K^{(1)}$

\[
T_p^\Omega=T_p:\mathcal M_K\longrightarrow\deleted{\unsure{\operatorname{Rep}(\pi,\mathbb Q_p)}}\,\operatorname{Mod\,f}(\mathbb Q_p)
\qquad\text{foncteur fibre en }\Omega
\]
\[
G^\Omega=G^{\deleted{\unsure{\Omega}}}=\operatorname{Aut}(T_p)
\qquad\text{pro-groupe alg. sur }\mathbb Q_p
\]
\[
\pi\xrightarrow{\rho}G(\mathbb Q_p)\qquad
\begin{aligned}
&G\text{ est l’enveloppe algébrique de }\pi\\
&(\mathrm{NB}\quad G/G^{\unsure{\circ|0}}\ \text{non tordu})
\end{aligned}
\]

Pour $M\in\operatorname{Ob}\mathcal M$,
\[
G_M=\operatorname{Im}\bigl(G\longrightarrow\underline{\operatorname{Aut}}_{\mathbb Q_p}(T_p(M))\bigr)\subset\underline{\operatorname{Aut}}_{\mathbb Q_p}(T_p(M))
\]
\[
\bigl(\pi_{\deleted{\unsure{M}}}\xrightarrow{\rho_M}G_M(\mathbb Q_p)\quad\ldotp\ldotp\ldotp\ldotp\ldotp\ \bigr)
\]

\diagram{\resizebox{\linewidth}{!}{$\begin{array}{lllll}
\substack{\text{structure}\\\text{sur }(\Omega,\pi)}
&\begin{tikzcd}[ampersand replacement=\&,row sep=0pt,column sep=3pt] \&G^{!}_{\Omega}\\G_{m_\Omega}\arrow[ur,"{\unsure{j|J}^{!}}"]\&\end{tikzcd}
&\substack{\unsure{\Phi|P|\mathfrak{P}}\\\text{avec section }e\\\text{sur }\Omega\\\text{\unsure{pas invariante|non invariante} sous }\pi}
&G_\Omega&\\[4pt]
\text{id}
&\begin{tikzcd}[ampersand replacement=\&,row sep=0pt,column sep=3pt] \&G^{\mathrm{Hdg}}_{\Omega}\\G_{m_\Omega}\arrow[ur,"{j_\Omega}"]\&\end{tikzcd}
&Q_\Omega=\unsure{\mathcal T|\tau}\mathbin{\wedge}_{G_{m_\Omega}}\Phi&G_\Omega&\\[4pt]
\substack{\text{par descente}\\\text{sur }K}
&\begin{tikzcd}[ampersand replacement=\&,row sep=0pt,column sep=3pt] \&G^{\mathrm{Hdg}}_{K}\\G_{m_K}\arrow[ur,"{j_K}"]\&\end{tikzcd}
& Q_K&G_K&T_{\mathrm{Hdg}}^{\deleted{\unsure{K}}}(M)=Q_K\mathbin{\wedge}_G M_K
\end{array}$}}
\margin{$\bigl[\mathcal T\deleted{\unsure{=}}\inserted{=\mathcal T^\Omega}$ $G_{m_\Omega}$-tor, sur $(\Omega,\pi)$ associé à $T_p(\Omega^\times)\otimes_{\mathbb Z_p}\Omega$ $\bigr]$}

Pour $M\in\operatorname{Ob}\widetilde{\mathrm{BT}}_V$, on a \deleted{\unsure{sur} $T_{\mathrm{Hdg}}(M)$} un

a) Un $\deleted{\unsure{K}}L$-vectoriel \deleted{\unsure{filtré}} $T_{\unsure{\mathrm{dR}|\mathrm{DR}|\mathrm{cr}}}(M)$ \deleted{et un isom}

\deleted{$\operatorname{gr}T_{\mathrm{DR}}(M)_K\simeq T_{\mathrm{Hdg}}(M)$}
% archive sheet number: 22
